%% kapa: exam=1 round=28 year=2017 period=1 form=A subject=civil_law subject_ko=민법 seq=23 number=23
%% source: Q-Net 감정평가사 기출문제 (cst00309) artlSeq=5199603 file=extracted/1008/5199603/2017년 제28회 감정평가사 1차시험 문제지/1교시 A형.pdf
%% conversion: kordoc 4.24.51 -> tools/gen1.py (auto); review=proofread-sonnet; flags=none
\begin{question}{23}
등기의 추정력에 관한 설명으로 옳지 않은 것은? (다툼이 있으면 판례에 따름)
\begin{choices}
  \item 신축건물에 소유권보존등기가 된 경우, 그 명의자가 신축한 것이 아니라도 그 보존등기는 실체관계에 부합하는 유효한 등기로 추정된다.
  \item 소유권이전등기 명의자는 제3자뿐만 아니라 전(前)소유자에 대해서도 적법한 등기원인에 의하여 소유권을 취득한 것으로 추정된다.
  \item 소유권이전등기는 등기원인과 절차가 적법하게 마쳐진 것으로 추정된다.
  \item 종중재산에 대한 유효한 명의신탁의 경우, 등기의 추정력에도 불구하고 신탁자는 수탁자에 대하여 명의신탁에 의한 등기임을 주장할 수 있다.
  \item 소유권이전청구권 보전을 위한 가등기가 있다고 하여 소유권이전등기를 청구할 수 있는 법률관계가 존재한다고 추정되는 것은 아니다.
\end{choices}
\end{question}
\answer{1}
